
\chapter{Национальные карточные сети СНГ}\label{ch:regional-schemes}
Банку или поставщику платёжных услуг (\textit{payment service provider}, PSP)
при~выходе в~соседнюю страну недостаточно терминала с~EMV, хоста с~ISO~8583
и~знакомого логотипа на~карте. Логотип не~показывает, кто управляет внутренним
трафиком, где завершается расчёт, обязателен ли внутренний приём и~как
кобейджинг (выпуск карты с~двумя платёжными приложениями двух разных сетей;
разбор см.~в~разделе~\ref{sec:up-mir}) открывает внешний путь карты
за~пределами страны.

\section{Мотивация национальных сетей}\label{sec:regional-why}

\highlight{Суверенный контроль над~правилами сети и~непрерывностью операций}
требует маршрутизации внутреннего трафика силами собственного оператора.
Маршрутизация через внешнюю сеть ставит весь внутренний оборот в~зависимость
от~остановки, санкционного давления или единоличного решения её владельца.
Регулятору нужна инфраструктура, которая работает независимо от~решений внешних
участников.

Национальная сеть и~процессинг убирают зависимость от~чужой валютной и~расчётной
инфраструктуры. Расчёты в~национальной валюте между участниками не~требуют
прохождения через корреспондентскую сеть (систему счетов, которые банки держат
друг у~друга для~исполнения межбанковских переводов) или зарубежный клиринг.
Порог входа ниже, чем у~глобальной сети. Прямое членство (статус участника
  сети, который даёт прямое подключение к~сети и~обязанности по~правилам,
в~отличие от~косвенного участия через спонсора) в~глобальной сети требует
капитала, технической сертификации и~инфраструктуры сопровождения, которой
у~небольшого регионального банка может не~быть.

Мир создавался под~санкционным давлением. После того как Visa и~Mastercard
в~2014~году приостановили обслуживание ряда российских банков, основана
Национальная система платёжных карт (НСПК;
см.~гл.~\ref{ch:mir})~\cite{nspk-about-company} и~в~2015~году запущена карта
Мир. В~марте 2022~года Visa и~Mastercard полностью прекратили операции
в~России, и~Мир остался основной внутренней карточной сетью. Иранская сеть
Shetab появилась раньше, в~2002~году, как проект унификации фрагментированной
карточной инфраструктуры страны. Санкции отрезали Иран от~международных сетей
позднее, превратив Shetab в~единственную доступную внутреннюю инфраструктуру.
UnionPay в~Китае выросла из~задачи обслуживать национальный розничный рынок~---
задолго до~того, как международные карты заняли там значимую долю.

Транзакции через собственную сеть генерируют interchange и~процессинговые
комиссии, которые остаются в~национальной экономике, а~не~перечисляются
иностранному оператору. Тарифы и~правила операций устанавливает оператор
расчётной системы внутри страны.

С~2024~года к~этому добавился санкционный механизм: включение оператора чужой
сети в~список SDN переносит риск на~эмитентов в~стране-партнёре. Банк, выпустивший карту
сети~X, после санкций против~X выбирает между нарушением правил сети и~риском
для~корреспондентских отношений в~валютах третьих стран; разбор
см.~в~разделе~\ref{sec:regional-after-sdn}. \highlight{Собственная сеть
устраняет этот риск по~устройству}: суверен сам определяет применение
санкций внутри страны и~не~подчинён внешним спискам SDN.

Белкарт, HUMO и~ЭЛКАРТ имеют собственный стандарт карточного приложения,
собственные данные авторизационного сообщения и~собственные правила отклонения.
Совместимость со~стеком EMV у~национальной сети не~следует из~самого факта
карточной эмиссии: перед интеграцией проверяется документация национальной сети,
а~не~общий EMV как точка отсчёта.

\section{Критерии сравнения}\label{sec:regional-framework}

Инженерное сравнение национальной сети начинается с~пяти параметров: оператор
и~физическое размещение процессинга; обязательность внутреннего приёма
(\textit{domestic acceptance}); устройство расчётной системы и~условия участия;
необходимость кобейджинга для~внешнего использования; степень совпадения
с~привычным стеком EMV/ISO~8583 из~глав~\ref{ch:emv} и~\ref{ch:tx-lifecycle}.

Формат сообщения и~чип ничего не~говорят о~географии приёма, обязательности
эквайринга и~распределении рисков работы сети. Различия проявляются
в~операторской модели~--- кто владеет сетью и~выпускает правила; в~размещении
процессинга и~расчёта~--- внутренняя обработка идёт через собственный узел,
и~тот же~оператор завершает урегулирование; в~модели приёма~--- обязательная
внутри страны или добровольная; во~внешнем пути карты~--- самостоятельном или
открываемом через кобейджинг.

Пока неизвестно, где находится расчётная система, как оформляется участие
и~обязателен ли внутренний приём, решение о~поддержке сети преждевременно:
главные ограничения коренятся в~организационной модели, а~транспортный формат
сообщения вторичен.

\section{Белкарт и~ArCa}\label{sec:regional-belkart-arca}

Белкарт~--- национальная сеть с~собственной внутренней инфраструктурой,
оформленной как бренд и~набор правил для~участников. Оператор системы~--- ОАО
\enquote{Банковский процессинговый центр}; правила системы охватывают эмиссию,
эквайринг, процессинг и~расчёты по~операциям с~карточками
Белкарт~\cite{belkart-about}.

До~2022~года среди продуктов Белкарт была кобейджинговая карта
Белкарт--Maestro (Maestro~--- дебетовый бренд Mastercard, закрытый
к~2023~году) для~операций в~Беларуси и~за~рубежом~\cite{belkart-cards};
кобейдж Белкарт--Мир был запущен в~2020~году, а~после приостановки Mastercard
в~регионе стал доминирующим внешним маршрутом.

ArCa с~самого начала работает на~инфраструктуре процессора, который обслуживает
и~международные сети. Armenian Card описывает себя как оператора
и~процессинговый центр национальной платёжной системы ArCa: её инфраструктура
обрабатывает транзакции и~проводит клиринг, а~банки-участники, помимо
Mastercard, Visa, American Express, DCI (Diners Club International), Мир и~JCB
(\textit{Japan Credit Bureau}~--- японская международная карточная сеть),
выпускают и~обслуживают карты ArCa~\cite{arca-about}.

\highlight{Процессинг разделяет логотип на~карте и~инфраструктуру расчёта}.
Даже компактные сети вроде Белкарт и~ArCa требуют
отдельного проекта, если банк встраивает их в~процессинг с~поддержкой
нескольких сетей: операторская модель и~обязательный внутренний приём задают
требования, специфичные для~конкретной страны, а~группа, работающая сразу
в~нескольких юрисдикциях региона, ведёт отдельный набор правил подключения
и~сопровождения по~каждой сети.

\section{HUMO, UzCard и~ЭЛКАРТ}\label{sec:regional-humo-uzcard-elkart}

Официальный FAQ HUMO называет HUMO национальной платёжной картой и~связывает
появление системы с~президентским указом о~создании National Interbank
Processing Centre (NIPC, Узбекистан~--- отдельная организация, не~путать
с~одноимённым по~смыслу IPC в~Кыргызстане, см.~ниже). Международные карты
внутри Узбекистана могут обрабатываться через HUMO в~роли второго процессора,
а~за~рубежом принимаются в~сети той системы, с~которой карта была
выпущена~\cite{humo-faq}.

UzCard на~своём официальном сайте называет себя ведущим национальным
платёжным оператором Республики Узбекистан. В~сообщении 2025~года UzCard
и~UnionPay объявили о~планах расширять выпуск кобейдж-карт и~обеспечивать
совместимость QR-кодов между сетями~\cite{uzcard-unionpay-2025}. Интеграция
банкоматных сетей HUMO и~UzCard прошла летом 2023~года: после президентского
указа от~9~марта 2023~года к~концу июля сети объединили. В~ноябрьском сообщении
HUMO сказано, что число банкоматов, доступных населению, выросло
вдвое~\cite{humo-uzcard-atm-2023}. Пара HUMO и~UzCard задаёт интегратору
отдельную задачу интероперабельности: поддержать обе карты и~отделить участки,
где инфраструктура уже объединена, от~участков, где правила операторов
различаются.

Официальный сайт ЭЛКАРТ выделяет среди базовых продуктов \enquote{Социальные
выплаты} и~\enquote{Пенсионную карту}; описание пенсионной карты прямо
связывает её с~получением пенсий, пособий и~компенсаций~\cite{elkart-home}.
В~том же FAQ сказано, что ЭЛКАРТ~--- национальная платёжная карточная система
под~управлением Interbank Processing Center (IPC~--- оператор национальной
  платёжной системы Кыргызстана \enquote{ЭЛКАРТ}; не~путать с~узбекским NIPC
выше), а~в~роли эмитента участвует и~Центральное казначейство Министерства
финансов Кыргызстана~\cite{elkart-home}. Страница IPC подтверждает, что
компания~--- оператор национальной платёжной системы \enquote{ЭЛКАРТ}
и~одновременно оказывает услуги по~обработке платежей и~проведению расчётов
в~международных платёжных системах, действующих на~рынке
страны~\cite{ipc-elkart-company}. Помимо кобейджа с~UnionPay у~ЭЛКАРТ
существует и~кобейдж с~Мир~--- продукт \enquote{ELCARD MIR NFC}, выпускаемый,
в~частности, Кыргызско-Швейцарским банком. Карта принимается в~обеих сетях.

HUMO использует собственные правила эмиссии и~расчёта; UzCard делит
национальный рынок с~HUMO и~подключается к~ней в~части банкоматной сети;
ЭЛКАРТ опирается на~государственные и~социальные сценарии. Выбор локальной
сети задаёт интегратору параметризацию терминала, маршрутизацию на~стороне
хоста, подключение ТСП, собственный список проверок релиза и~отдельную
поддержку операций.

\section{Азербайджан: AzeriCard и~MilliKart}\label{sec:regional-azerbaijan}

\highlight{В~Азербайджане локальная инфраструктура сосредоточена в~процессинговых центрах без~собственного карточного бренда}. AzeriCard
и~MilliKart~--- два независимых процессинговых центра, ни~один из~них
не~эмитирует карту собственного бренда.

AzeriCard работает с~1997~года при~International Bank of Azerbaijan,
обслуживает более~30 банков страны и~сертифицирован как принципал Visa,
Mastercard, American Express, Diners Club, UnionPay
и~JCB~\cite{azericard-overview}. MilliKart основан Центральным банком
Азербайджана в~2006~году, с~октября того же~года работает как ООО
с~17~банками-участниками и~ЕБРР в~капитале и~имеет лицензии полного членства
(\textit{full membership}) Mastercard, Visa
и~UnionPay~\cite{millikart-overview}.

Клиент держит карту с~международным логотипом; внутренняя транзакция проходит
через азербайджанский процессинговый центр и~не~передаётся зарубежному сервису;
внешний путь карты работает по~правилам Visa, Mastercard или~UnionPay
без~отдельного кобейджинга. Для~интегратора это снижает требования к~хосту~---
сертификация ISO~8583-сообщения у~процессора уже покрывает международные
сети~--- и~одновременно повышает зависимость от~их правил: своей карточной
сети, к~которой можно было бы переключиться в~случае санкционного давления,
у~такой модели нет.

В~2024~году AzeriCard прошёл интеграцию с~Apple Pay и~Google Pay через шлюз
Akurateco, обеспечив национальное покрытие NFC-кошельков. Процессор
не~эмитирует карты, но контролирует путь токенизации: сертификация в~Visa Token
Service и~Mastercard Digital Enablement Service ведётся от~его имени
для~банков-участников. Ту~же модель MilliKart и~AzeriCard применяют
к~внутренней QR-инфраструктуре, где локальный процессор задаёт стандарт
QR-кода, а~банки подключаются как участники, а~не~как операторы.

\section{Казахстан без~национальной карточной сети}\label{sec:regional-kazakhstan}

\highlight{В~Казахстане место национальной сети заняла частная платёжная
экосистема Kaspi.kz}. По~собственной отчётности Kaspi обрабатывает
около~80\,\% розничных платёжных операций в~стране и~обслуживает
14~млн ежемесячно активных
пользователей при~населении 20~млн~\cite{kaspi-payments-platform}. На~июль
2024~года в~Казахстане в~обращении 78,1~млн платёжных карт; 80\,\% дебетовые,
15,9\,\% кредитные, 4,1\,\% предоплаченные (\textit{prepaid}) и~дебетовые
с~лимитом~\cite{nbk-cards-statistics}. Все они выпущены под~брендами Visa,
Mastercard, UnionPay и~Мир~--- собственного карточного бренда страна
не~создала.

Нацбанк~РК в~2024~году объявил о~запуске универсальной межбанковской
QR-инфраструктуры с~приёмом в~любом банке страны. В~ответ Kaspi в~сентябре
2024~года открыл собственную платёжную сеть для~клиентов Home Credit Bank
и~Altyn Bank, до~этого работавших только с~картами своих
эмитентов~\cite{kaspi-opens-payments-2024}. Аргументация повторяет логику
национальной сети: контроль над~внутренней маршрутизацией, независимость
от~международных операторов, собственные правила приёма. Разница в~субъекте:
задачу решает частный игрок, а~не~государственный оператор расчётной системы.

Для~интегратора, выходящего в~Казахстан, выбор состоит из~двух подключений:
Kaspi через её API либо банки через QR-сервис Нацбанка. Доля Kaspi не~позволяет
игнорировать её; универсальный QR-сервис Нацбанка пока охватывает меньший
поток, но обеспечивает регуляторное равенство участников. Платёжный суверенитет
здесь обеспечивает не~операторская модель национальной сети, а~монополия
местной частной экосистемы, согласовавшей с~регулятором правила доступа третьих
сторон.

Когда доля единственного провайдера превышает половину розничных платежей,
регулятор сталкивается с~теми же задачами, что и~центральный банк страны
без~собственной сети: контроль непрерывности операций, санкционная
политика, тарификация. Разница в~субъекте отчётности: в~национальной сети
оператор подотчётен регулятору, в~частной экосистеме~--- акционерам и~биржевому
рынку. Регуляторный ответ Нацбанка~РК (универсальная QR-инфраструктура,
межбанковский доступ) повторяет шаги, которые десятилетие назад привели
к~созданию НСПК и~Белкарта; стимулом здесь стал частный провайдер,
а~не~санкционное давление.

\section{Кобейджинг после внесения НСПК в~список SDN}\label{sec:regional-after-sdn}

23~февраля 2024~года Минфин США внёс АО~НСПК в~список SDN~\cite{ofac-jy2117-nspk}. Для~национальных сетей кобейдж
с~Мир из~канала внешнего использования превратился в~канал санкционной
экспозиции. Армянская~ArCa свернула продукт ArCa--Мир к~30~марта 2024~года
(см.~таблицу~\ref{tab:regional-comparison}); зарубежные эмитенты, у~которых
кобейдж с~Мир оставался активным, пересматривали набор карточных продуктов из-за
риска вторичных санкций для~их банков-корреспондентов.

К~2025~году Мир принимается, по~открытым заявлениям ЦБ~РФ, в~13~странах;
в~девяти~--- с~ограничениями (Армения, Венесуэла, Вьетнам, Казахстан, Лаос,
Молдова, Мьянма, Таджикистан, Никарагуа)~\cite{mir-acceptance-tass-2025}. В~мае
2025~года ЦБ Ирана объявил о~завершении второго этапа интеграции с~Мир,
обеспечив приём карт по~бесконтактной технологии в~местной сети эквайринга.
Список стран и~условия приёма пересматриваются чаще, чем выходит печатное
издание; интегратор сверяется с~актуальными ограничениями национального
оператора-партнёра, а~не~с~общим списком стран.

\highlight{После февраля 2024~года решение по~кобейджу начинается
с~санкционной оценки; продуктовая экономика идёт следом}. Для~зарубежного эмитента,
выпустившего карту X--Мир, вопрос о~поддержке продукта зависит от~оценки
риска по~OFAC и~требований банков-корреспондентов, а~не~от~экономики карточного
канала. Для~эмитента в~стране, где обязательный приём собственной национальной
сети требует кобейджа для~внешнего пути, остаются два варианта: разрабатывать
альтернативный маршрут через~UnionPay либо~JCB или отказываться от~внешнего
использования вовсе. Третий вариант~--- оставлять кобейдж с~Мир активным~--- после
февраля 2024~года теряет смысл для~банков, сохраняющих корреспондентские
отношения в~долларах и~евро.

Кобейдж с~UnionPay сохраняет внешнее использование в~Китае, странах ASEAN,
на~Ближнем Востоке и~в~ряде европейских юрисдикций, но требует сертификации
по~правилам UnionPay International и~подключения к~CIPS (\textit{Cross-Border
Interbank Payment System}) для~межбанковского клиринга в~юанях. Кобейдж с~JCB
даёт приём в~Японии и~Юго-Восточной Азии, но эквайринговая сеть там заметно
уже, чем у~UnionPay. Национальная сеть, обязательная к~приёму внутри страны,
продолжает работать независимо от~внешнего кобейджа: внутренний поток
к~оператору сети санкционная экспозиция не~затрагивает.

Параллельная инфраструктура~--- СПФС Банка России для~межбанковских сообщений
в~рублях и~национальных валютах стран-партнёров~--- к~2025~году расширилась
до~зарубежных участников из~стран~СНГ, Ирана и~ряда юрисдикций Ближнего Востока
и~Юго-Восточной Азии~\cite{cbr-spfs}. СПФС не~заменяет карточную маршрутизацию и~не~подменяет
Visa или~Mastercard в~розничном приёме, но обеспечивает
корреспондентский канал для~межбанковского клиринга, независимый от~SWIFT
и~доступный участникам в~условиях санкционных ограничений. Свёртывание кобейджа с~Мир
поэтому не~означает разрыв всех расчётных связей с~российскими банками
для~операторов национальных сетей СНГ: рублёвый и~национальный поток продолжается
через СПФС, а~не~через карточную инфраструктуру.

\section{Последствия для~интегратора}\label{sec:regional-integrator}

Технической совместимости с~EMV из~главы~\ref{ch:emv} и~знакомого карточного
жизненного цикла из~главы~\ref{ch:tx-lifecycle} банку, процессингу и~шлюзу
недостаточно: эти стандарты не~отвечают на~вопросы, кто задаёт правила сети,
где обязателен приём, как устроен кобейджинг и~насколько широк охват внешнего
использования.

Продавец сталкивается с~той~же асимметрией: в~одной стране национальная сеть
обязательна для~приёма, в~другой это дополнительный внутренний канал, важный
для~отдельных категорий клиентов или выплат. От~этого зависят витрина оплаты,
приоритет брендов в~интерфейсе, формулировки приёма в~договоре и~запасные
сценарии службы поддержки.

В~сводной таблице~\ref{tab:regional-comparison} BIN-диапазоны (\textit{bank
  identification number}; в~ISO/IEC~7812~--- стандарте нумерации идентификаторов
эмитентов~--- обозначается как \textit{issuer identification number}, IIN)
для~всех сетей, кроме~Мир, требуют сверки по~реестру IIN
ISO/IEC~7812~\cite{iso-7812-transition} и~публичной документации операторов:
в~открытых источниках встречаются устаревшие значения локального внутреннего
распределения, не~зарегистрированные в~ISO.

\begin{table}[t]
  \caption{Параметры национальных карточных сетей
  СНГ.}\label{tab:regional-comparison}
  \begingroup
  \tablesetup
  \setlength{\tabcolsep}{3pt}
  \begin{tabularx}{\textwidth}{
      @{}P{0.14\textwidth}
      Y
      Y
      Y
      Y@{}
    }
    \toprule
    \headrow
    \thd{Сеть} & \thd{Диапазон BIN}
    & \thd{Внутренний interchange}
    & \thd{Обязательный приём} & \thd{Внешний путь} \\
    \midrule
    Мир & 2200--2204 & По~таблицам НСПК
    & Да (ТСП свыше~20~млн руб./год)
    & Кобейдж с~UnionPay \\
    \rowsep
    Белкарт & требует сверки & По~правилам оператора
    & Да (по~ряду категорий ТСП)
    & Кобейдж с~Мир; до~2022~--- Maestro \\
    \rowsep
    HUMO & требует сверки & По~правилам НСПК Узбекистана
    & Да (внутри Узбекистана)
    & Кобейдж с~UnionPay \\
    \rowsep
    UzCard & требует сверки & По~правилам UzCard
    & Да (внутри Узбекистана)
    & Ограниченный \\
    \rowsep
    ArCa & требует сверки & По~правилам ЦБ Армении
    & Да (внутри Армении)
    & Кобейдж с~Visa/MC, JCB, UnionPay; ArCa--Мир~--- до~30~марта 2024 \\
    \rowsep
    ЭЛКАРТ & требует сверки & По~правилам оператора
    & Да (внутри Кыргызстана)
    & Кобейдж с~UnionPay и~Мир \\
    \bottomrule
  \end{tabularx}\par
  \endgroup
\end{table}

Поддержка национальной карты~--- это согласие на~её внутреннюю платёжную систему
целиком: операторскую модель, требования сертификации, обязательность приёма
и~санкционную экспозицию кобейджа.

\practicenote{%
  Минимальный список проверок для~интегратора по~национальной сети:
  где находится оператор и~расчётная система, обязателен ли приём
  внутри страны, нужен ли отдельный кобейджинговый путь для~внешнего
  использования, какие сертификации терминала и~хоста требуются
  и~кто отвечает за~международный приём. Если хотя бы один пункт
  остаётся неясным, интеграцию лучше считать незавершённой, даже
  если лабораторный платёж по~карте уже проходит.
}

Банковская группа с~дочерними банками в~нескольких странах СНГ ведёт пять или
шесть параллельных интеграций с~национальными сетями. У~Белкарт, HUMO, UzCard,
ArCa и~ЭЛКАРТ собственное EMV-приложение (AID) и~собственное терминальное ядро:
параметры терминала, профиль EMV-приложения, структура авторизационного
сообщения и~правила отклонения у~каждой сети свои.
Расчётный день и~формат сводных файлов клиринга тоже различаются; единого расчётного
календаря у~группы не~будет. Унификация достигается собственной шиной
процессинга: один формат входящего сообщения и~маршрутизация по~диапазону BIN
на~подключение нужного оператора. То~же распределение действует
для~производства и~персонализации карт (раздел~\ref{sec:emv-personalization}): профиль приложения готовится отдельно
под~каждую сеть, а~скрипты подготовки данных и~договор с~производителем
карточных заготовок остаются общими. Экономика такой группы строится
на~отделении общей инфраструктуры процессинга от~локальной специфики каждой
сети; расчёт на~унификацию стандартов в~этой группе не~работает.
